\documentclass[11pt]{article}

\usepackage[margin=25mm]{geometry}

\usepackage{amsmath,amssymb,fontspec,unicode-math}

\IfFontExistsTF{DejaVu Serif}{\setmainfont{DejaVu Serif}}{\setmainfont{Latin Modern Roman}}

\IfFontExistsTF{Latin Modern Math}{\setmathfont{Latin Modern Math}}{\setmathfont{DejaVu Math TeX Gyre}}

\usepackage{xurl,hyperref}

\setlength{\parindent}{0pt}\setlength{\parskip}{6pt}

\begin{document}

\title{Learnable Burst-Encodable Time-of-Flight Imaging for High-Fidelity Long-Distance Depth Sensing}

\author{ELVIA preliminary AI analysis}\date{21 October 2026}\maketitle

Preliminary assessment: Minor.\par

Independent scientific review remains pending. Author adoption has not been requested.

The public forward model assumes the same ambient exposure for independently learned codes, and its implemented information expression differs from the Gaussian Fisher formula. No historical depth result is remeasured.

\section*{Review scope}

Latest-version, claim-focused review of the cited author-paper passages and applicable current author-release paths. The stated counterexamples and scope controls were checked analytically; no author program, old verifier, training run or benchmark was executed.

\section*{Problem 1: Ambient light is not integrated through each learned code}

Source: Section 3.1, PDF p.4, Eqs.(5)–(6); network.py, lines 131–178

Assessment: supported; impact: minor

\subsection*{Source claim (paraphrase)}

Equations (5)–(6) require ambient contribution proportional to \(I_{\rm amb}\int D_i(t)\,dt\). `network.py` instead adds a shared factor based on \(500\) to every \(1000\)-sample channel.

\subsection*{Analysis}

Independent normalization does not force equal code sums. For valid code sums \((250,750,500,200)\) and ambient scale \(0.6\), the coded contributions are \((150,450,300,120)\); the shared implementation gives \((300,300,300,300)\). Half-duty channels are a passing special case.

\subsection*{Finding and scope}

The public forward model imposes an unenforced equal-duty assumption. Its historical measurement effect is unknown.

\subsection*{Proposed correction}

Use the actual per-channel code sum, including sample spacing and accumulation factors, in both expected intensities and information-loss curves. Derive noise variance from those same corrected means.

An explicit local correction is described; author adoption and a corresponding published erratum are not established.

\section*{Problem 2: The information expression is not the stated Fisher information}

Source: Section 3.3, PDF p.6, Eq.(9); loss\_function.py, lines 6–23; train.py, line 186

Assessment: supported; impact: minor

\subsection*{Source claim (paraphrase)}

For Gaussian measurements with mean \(\mu(d)\) and variance \(V(d)=\mu(d)+\nu\), Eq.(9) requires \(J=\mu^{\prime2}/V+V^{\prime2}/(2V^2)\). `loss\_function.py` introduces a second-derivative term and different coefficients.

\subsection*{Analysis}

Taking the expectation of the squared Gaussian score yields the displayed nonnegative Fisher expression. A second derivative of the mean does not appear in that identity. A finite positive-variance quadratic mean can make the implemented sum negative while Eq.(9)’s finite-difference sum is positive, so the discrepancy is not a common normalization. Such a quadratic is a diagnostic mean, not an asserted learned code curve.

\subsection*{Finding and scope}

The routine does not compute Eq.(9). This does not prove an observed optimization failure in the paper.

\subsection*{Proposed correction}

Use \([1/V+1/(2V^2)](g\mu)^2\) under the same mean-dependent variance model. State finite-difference boundaries, the depth/phase Jacobian and the design distribution used for accumulation.

An explicit local correction is described; author adoption and a corresponding published erratum are not established.

\section*{Limits}

Preliminary AI assessment; independent human scientific adjudication remains pending.\par

An analytic counterexample or a release discrepancy does not by itself remeasure the paper’s reported experiments.\par

\section*{Source identity}

NeurIPS 2025 official proceedings PDF; compared with current arXiv:2505.22025v1

\url{https://proceedings.neurips.cc/paper_files/paper/2025/file/bc43bf4c31d5f074d023ab015a6107c4-Paper-Conference.pdf}

Reviewed PDF SHA-256: \nolinkurl{cb4e722addd104d648a89b71a03764f7abde0378fe71cd67f97d799f6759b2db}

The selected author/publisher source is the same work as the original case, matched by exact title, author identity and official version channel. The exact selected PDF SHA-256 was recomputed in the latest-version pass for Issue \#27. The PDF is a retained author/publisher copy; its original retrieval time is stated only when supported by matching acquisition metadata. File modification time was not substituted.

\end{document}
